% Bagian: second-order-logic
% Bab: syntax-and-semantics
% Subbagian: language-of-sol

\documentclass[../../../include/open-logic-section]{subfiles}

\begin{document}

\olfileid[id]{sol}{syn}{lan}

\olsection{Bahasa Logika Orde Kedua}

Seperti dalam logika orde pertama, ekspresi logika orde kedua dibangun dari
suatu kosakata dasar yang memuat \emph{!!{variable}s},
\emph{!!{constant}s}, \emph{!!{predicate}s}, dan kadang-kadang
\emph{!!{function}s}. Dari unsur-unsur tersebut, bersama penghubung logis,
kuantor, dan tanda baca seperti tanda kurung dan koma, dibentuk \emph{suku}
dan \emph{!!{formula}s}. Perbedaannya ialah bahwa, selain variabel untuk
objek, logika orde kedua juga memuat variabel untuk relasi dan fungsi, serta
memperbolehkan kuantifikasi atas variabel-variabel itu. Jadi, simbol logis
logika orde kedua ialah simbol logis logika orde pertama, ditambah:

\begin{enumerate}
\item Suatu himpunan !!{denumerable}s !!{variable}s relasi orde kedua bagi setiap
  aritas~$n$: $\Obj V_0^n$, $\Obj V_1^n$, $\Obj V_2^n$, \dots
\item Suatu himpunan !!{denumerable}s !!{variable}s fungsi orde kedua:
  $\Obj u_0^n$, $\Obj u_1^n$, $\Obj u_2^n$, \dots
\end{enumerate}

Dalam logika orde pertama, !!{identity}~$\eq$ biasanya disertakan. Dalam
logika orde pertama, simbol-simbol nonlogis suatu bahasa~$\Lang{L}$ sangat
penting agar kita dapat menyatakan sesuatu yang menarik. Tentu ada
!!{sentence}s yang tidak menggunakan simbol nonlogis, tetapi dengan hanya
$\eq$, sulit untuk mengatakan sesuatu yang menarik. Dalam logika orde kedua,
karena kita mempunyai persediaan variabel relasi dan fungsi yang tidak
terbatas, kita dapat mengatakan segala sesuatu yang dapat kita katakan dalam
suatu bahasa orde pertama, bahkan tanpa persediaan khusus simbol nonlogis.

\end{document}
